\song{Grónská písnička (Jaromír Nohavica)}
%Mara_2016_12_25

\ve{}
\ch{D}Daleko na \ch{Emi}severu je \ch{A7}Grónská \ch{D}zem,\\
\ch{D}{ž}ije tam \ch{Emi}Eskymačka s \ch{A7}Eskymá\ch{D}kem,\\
\ch{D}my bychom \ch{Emi}umrzli, jim \ch{G}není zima, \ch{D}\\
\ch{D}snídají \ch{Emi}nanuky \ch{A7}a eskima. \ch{D}\\
\ch{D}my bychom \ch{Emi}umrzli, jim \ch{G}není zima, \ch{D}\\
\ch{D}snídají \ch{Emi}nanuky \ch{A7}a eskima. \ch{D}

\ve{}
Mají se bezvadně, vyspí se moc,\\
půl roku trvá tam polární noc,\\
na jaře vzbudí se a vyběhnou ven,\\
půl roku trvá tam polární den.

\ve{}
Když sněhu napadne nad kotníky,\\
hrávají s medvědy na četníky,\\
medvědi těžko jsou k poražení,\\
neboť medvědy ve sněhu vidět není.

\ve{}
Pokaždé ve středu, přesně ve dvě\\
zaklepe na na íglů hlavní medvěd:\\
\uv{Dobrý den, mohu dál na vteřinu?\\
Nesu vám trochu ryb na svačinu.}

\ve{}
V kotlíku bublá čaj, kamna hřejí,\\
psi venku hlídají před zloději,\\
smíchem se otřásá celé íglů,\\
medvěd jim předvádí spoustu fíglů.

\ve{}
Tak žijou vesele na severu,\\
srandu si dělají z teploměrů,\\
my bychom umrzli, jim není zima,\\ 
neboť jsou doma a mezi svýma.

\esong